% from ega4-18-fr.tex
\oldpage[IV-4]{158}

est plat (2.3.4); d'autre part, $\eta$ est aussi point générique de $g^{-1}(y_1)$ ($0_{\mathrm I}$, 2.1.8), et on sait par ailleurs que les fibres de $g$ sont des espaces discrets (18.8.12), donc $g^{-1}(y_1)=\{\eta\}$, et par suite $y'_1=\eta$. Le morphisme $\Spec(\sh{O}_{X',x'})\to\Spec(\sh{O}_{Y',y'})$ restriction de $f'$ est donc dominant, et puisque $\sh{O}_{Y',y'}=A'$ est intègre, l'homomorphisme correspondant $\sh{O}_{Y',y'}\to\sh{O}_{X',x'}$ est injectif (I, 1.2.7). D'autre part, $f'$ est non ramifié au point $x'$, (donc (18.8.3, $d$)), il y a un voisinage $U$ de $x'$ dans $X'$ tel que $f'|U$ soit une immersion fermée de $U$ dans $Y'$; \emph{a fortiori} l'homomorphisme $\sh{O}_{Y',y'}\to\sh{O}_{X',x'}$ est surjectif, et puisqu'il est injectif, il est bijectif. Mais alors $\sh{O}_{X',x'}$ est un $\sh{O}_{Y',y'}$-module plat, et $f'$ est donc étale au point $x'$ (17.6.1).

Enfin, lorsque $\sh{O}_{Y,y}$ est noethérien, on sait que si $f$ est étale au point $x$, on a $\dim\sh{O}_{Y,y}=\dim\sh{O}_{X,x}$ (17.6.4). Réciproquement, supposons vérifiées les conditions (i) et (ii \emph{bis}), et montrons qu'elles entraînent (ii). Soit $\mathfrak J$ le noyau de l'homomorphisme canonique $\sh{O}_{Y,y}\to\sh{O}_{X,x}$; restreignant au besoin $Y$ à un voisinage de $y$, on peut supposer que $\mathfrak J=\mathscr{J}_y$, où $\mathscr{J}$ est un Idéal de $\sh{O}_Y$; si $Y_1$ est le sous-préschéma fermé de $Y$ défini par $\mathscr{J}$, on peut, en restreignant encore au besoin $X$ et $Y$ à des voisinages de $x$ et $y$ respectivement, supposer que $f$ se factorise en $X\xrightarrow{f_1}Y_1\to Y$ (I, 6.5.1). Il est clair que $f_1$ est encore non ramifié au point $x$, donc (17.4.1) quasi-fini en ce point; la démonstration de (5.4.1, (i)) prouve alors que l'on a $\dim(\sh{O}_{X,x})\leq\dim(\sh{O}_{Y,y}/\mathfrak J)$. Mais si l'on avait $\mathfrak J\neq0$, on en conclurait que $\dim(\sh{O}_{Y,y}/\mathfrak J)<\dim(\sh{O}_{Y,y})$ puisque $Y$ est intègre ($0$, 16.1.2.2); on aurait par suite $\dim\sh{O}_{X,x}<\dim\sh{O}_{Y,y}$ contrairement à l'hypothèse, ce qui prouve (ii).

\begin{remark}[18.10.2]
\label{IV.18.10.2}
\begin{enumerate}
  \item[(i)] Lorsque $A=\sh{O}_{Y,y}$ est noethérien, et que l'on sait que son complété $\widehat A$ est intègre (par exemple si $A$ est régulier), on peut, dans la démonstration précédente, remplacer $A'$ par $\widehat A$.
  \item[(ii)] Lorsque $Y$ est localement noethérien, on peut donner de (18.10.1) une démonstration plus rapide, n'utilisant pas la hensélisation stricte, mais faisant intervenir les résultats assez délicats des \textsection\textsection~14 et 15. Comme $f$ est quasi-fini au point $x$ par hypothèse (17.4.1), on peut, en remplaçant $X$ par un voisinage ouvert de $x$, supposer que $f^{-1}(y)=\{x\}$. L'hypothèse (ii) entraîne, comme on l'a vu, l'existence d'une composante irréductible $Z$ de $X$ contenant $x$ et dominant l'unique composante irréductible $Y_0$ de $Y$ contenant $y$. Le morphisme $f|Z$ étant quasi-fini au point $x$, donc équidimensionnel en ce point (13.2.2), il résulte de l'hypothèse sur $Y$ et du critère de Chevalley (14.4.4) que $g=f|Z$ est universellement ouvert au point $x$. Par ailleurs puisque $f$ (donc aussi $g$) est non ramifié au point $x$, $g^{-1}(y)$ est géométriquement réduit sur $k(y)$ (17.4.1); comme $\sh{O}_{Y,y}$ est intègre, on conclut de (15.2.3) que $g$ est plat au point $x$, donc $f$ est étale en ce point (18.4.9).
  \item[(iii)] Les notations et les hypothèses sur $Y$ étant celles de (18.10.1), supposons seulement que $f$ soit localement de type fini (et non nécessairement localement de présentation finie). Alors, pour que $\sh{O}_{X,x}$ soit une $\sh{O}_{Y,y}$-algèbre essentiellement étale (18.6.1), il faut et il suffit que $f$ vérifie les deux conditions suivantes:

